% ====================================================================
% Control-plane / computational-cost limitation. Preventive honesty,
% same spirit as the K=4 audit. Addresses the serious shared critique
% about computational cost. Goes into the discussion/limitations section.
% ====================================================================
\subsection{The control-plane cost we are not hiding}
\label{sec:control-plane}

The reference core fits in about sixty lines of Python, and it is tempting to
read that as computational frugality. It is not, and we want to be explicit
about the gap, because it is the most consequential limitation of this work.

Two distinct notions of parsimony must be separated. The core is parsimonious in
\emph{mechanism}: two terms, no per-flow state, no precomputed path tables, no
hand-tuned controller. It is \emph{not} demonstrated to be parsimonious in
\emph{computation or control overhead}. Algorithm~\ref{alg:core} evaluates a
min-potential path by dynamic programming over the hop budget, $O(k|E|)$ per
decision, and to evaluate $\phi(u,v)$ on each edge it needs the instantaneous
load $x$ and scar state $s$ of every edge in the graph. Our flow-level simulator
provides that global state for free and instantaneously. A real deployment would
not: the load and loss state of remote edges would have to reach the deciding
node through telemetry or flooding, with latency and overhead that our model
omits entirely.

This matters for the comparison as much as for deployability. DRILL, in
particular, is a switch-local, per-hop mechanism that reacts to local queue
occupancy and assumes no global view; CONGA distributes congestion state but
within a bounded fabric. By granting our core an idealized, zero-latency global
state, the simulator may flatter it relative to a genuinely distributed,
local-information baseline. How the equivalence degrades when the global state is stale was, until this revision, the single most important experiment this work left open; Section~\ref{sec:freshness} now reports it, together with a per-decision cost benchmark and a bursty-arrival stress. We therefore frame the core's standing as
\emph{equivalence in routing quality under idealized state visibility}, not as a
claim of operational efficiency. Establishing whether the core survives the remaining constraints---telemetry bandwidth and propagation dynamics beyond the staleness and visibility priced in Section~\ref{sec:freshness}---is still required before any deployability claim could be made.
